% !TeX encoding = utf-8
% !TeX program = lualatex
% !TeX spellcheck = de_DE
% !BIB program = biber
\documentclass[%singlesided,
               doublesided,
               paper=a4,
               fontsize=10pt
              ]{my-resume}

\usepackage[ngerman]{babel}
\usepackage[
    % sorting=nyt,      % nyt | nty | nyt | nyvt | anyt | ynt | ydnt | none | debug
    % style=authoryear, % numeric | alphabetic | authoryear | authortitle | verbose | draft | debug
    % uniquename=false, % true, false , init, full, allinit, allfull, mininit, minfull
]{biblatex}
\addbibresource{Literature.bib}
\usepackage{csquotes}

%%%%%%%%%%%%%%%%%%%%%%%%%%%%%%%%%%%%%%%%%%%%%%%%%%%%%%%%%%%%%%%%%%%%%%%%%%%%%%%%
% set geometry
%%%%%%%%%%%%%%%%%%%%%%%%%%%%%%%%%%%%%%%%%%%%%%%%%%%%%%%%%%%%%%%%%%%%%%%%%%%%%%%%

\setlength\highlightwidth{8cm}
\setlength\headerheight{3.0cm}            % note that margintop gets added to this value, i.e. the header bar is 5cm
\setlength\marginleft{1cm}
\setlength\marginright{\marginleft}      % needs to be 1.5 times to be actually equal. why?
\setlength\margintop{1cm}
\setlength\marginbottom{1cm}


%%%%%%%%%%%%%%%%%%%%%%%%%%%%%%%%%%%%%%%%%%%%%%%%%%%%%%%%%%%%%%%%%%%%%%%%%%%%%%%%
% FONTS
%%%%%%%%%%%%%%%%%%%%%%%%%%%%%%%%%%%%%%%%%%%%%%%%%%%%%%%%%%%%%%%%%%%%%%%%%%%%%%%%

\RequirePackage{fontspec}
\setmainfont{Carlito}


%%%%%%%%%%%%%%%%%%%%%%%%%%%%%%%%%%%%%%%%%%%%%%%%%%%%%%%%%%%%%%%%%%%%%%%%%%%%%%%%
% COLORS
%%%%%%%%%%%%%%%%%%%%%%%%%%%%%%%%%%%%%%%%%%%%%%%%%%%%%%%%%%%%%%%%%%%%%%%%%%%%%%%%

\colorlet{highlightbarcolor}{lightgray}
\colorlet{headerbarcolor}{darkgray}

\colorlet{headerfontcolor}{white}
\colorlet{accent}{awesome-red}
\colorlet{heading}{black}
\colorlet{emphasis}{black}
\colorlet{body}{black}


%%%%%%%%%%%%%%%%%%%%%%%%%%%%%%%%%%%%%%%%%%%%%%%%%%%%%%%%%%%%%%%%%%%%%%%%%%%%%%%%
% set document
%%%%%%%%%%%%%%%%%%%%%%%%%%%%%%%%%%%%%%%%%%%%%%%%%%%%%%%%%%%%%%%%%%%%%%%%%%%%%%%%


\begin{document}

\name{Dr. Benedikt Kleinmeier}
\tagline{Informatiker / Leidenschaft für gute Software /  Interesse an Hardware und Physik}
\photo[round]{Figures/Portrait/bk}{\dimexpr \headerheight-\marginbottom}   % make photo exactly match the header with margintop/marginright/marginbottom as margin

\makeheader

\highlightbar{

    \section{Kontakt}
    
    \begin{infolist}
        \email{kleinmeier.benedikt@gmail.com}
        % \phone{}
        % \location{}
        \homepage{Projekte}{https://schwefelsaeure.github.io/portfolio/}
        \github{Github}{https://github.com/Schwefelsaeure/}
        % \linkedin{Name Surname}{https://www.linkedin.com/in/name-surname/}
        \orcid{Publikationen}{https://orcid.org/0000-0003-2344-8265}
        % \ads{}}
    \end{infolist}
        
    \section{Kenntnisse}
    
    \skillsection{Programmiersprachen}
    \skill{Assembler (x86, TriCore)}{2}
    \skill{Bash}{4}
    \skill{C}{5}
    \skill{C++}{3}
    \skill{C\#}{3}
    \skill{Java}{5}
    \skill{JavaScript}{2}
    \skill{Python}{5}
    \vspace{0.5em}
    
    \skillsection{Auszeichnungssprachen}
    \skill{HTML/CSS}{4}
    \skill{LaTeX}{4}
    \vspace{0.5em}
    
    \skillsection{Betriebssysteme}
    \skill{Linux}{5}
    \skill{MacOS}{1}
    \skill{Windows}{4}
    \vspace{0.5em}
    
    \skillsection{Software \& Tools}
    \skill{Visualisation}{3}
    (e.g. matplotlib, gnuplot, ...)\\
    \skill{Data handling/analysis}{3}
    (e.g. numpy, pandas, ...)\\
    \skill{Docker}{2}
    \skill{Office}{4}
    \vspace{0.5em}
    
    \skillsection{Vorgehensmodelle}
    ISO26262, Scrum, V-Modell
    \vspace{0.5em}
    
    \skillsection{Software-Methodiken}
    Continuous Integration, Test-Driven Devel"=opment
    \vspace{0.5em}
    
    \skillsection{Sprachen}
    \skill{Deutsch}{5}
    \skill{Englisch}{3}

      \section{Auszeichnungen}

    \cvachievement{\faTrophy}{2019}{TUM Graduate School Internationalization Grant}
    \cvachievement{\faTrophy}{2010}{Nominierung Stipendiatenprogramm I.C.S.}
    \cvachievement{\faTrophy}{2004}{Jahrgangsbester Abschlussjahrgang Realschule}

}
\mainbar{
    
    \section[\faGears]{Arbeitserfahrung}
    \job{seit 07/2021}
    {Infineon Technologies AG, Neubiberg}
    {Principal Software Engineer/Architect (In-House- und Kunden-Tools)}
    {\begin{itemize}
    		\item SDK- und GUI-Entwicklung (PyQt) von Radar-Applikationen zur Konfiguration von Radar-Sensoren, Daten-Akquise und Visualisierung.
            \item Weiterentwicklung einer Eclipse-basierten Applikation zur Ansteuerung von Laborequipment und zur Verwaltung von kryptografischen Keys.
            \item Entwicklung einer Workflow-Engine inklusive Web-Dashboard zur automatischen Prozessierung von Bestellungen.
            \item Entwicklung einer GUI-Applikation (Swing) zur Ansteuerung von SPICE-Simulatoren und Auswertung von Simulationsergebnissen.
            \item[] \tag{C++} \tag{Java} \tag{Python} \tag{GUI} \tag{PyQt} \tag{Eclipse/SWT} \tag{plotly/dash} 
    \end{itemize}}

    \job{12/2017 -- 05/2021}
        {Hochschule München / TU München / University of Sussex, München / Brighton}
        {Promotion}
        {\begin{itemize}
            \item Modellierung und Simulation von Verhaltensänderungen von Fußgängern.
            \item Umsetzung der Modellierung und Simulation im eigenen Open-Source-Simulator Vadere (\url{www.vadere.org}).
            \item GUI-Programmierung in Java sowie Datenauswertung mittels Python mit Fokus auf Test-Driven-Development und Continuous Integration.
            \item[] \tag{Java} \tag{Python} \tag{JUnit} \tag{pandas} \tag{matplotlib} \tag{GUI}
        \end{itemize}}
    
    \job{05/2015 -- 11/2017}
        {Infineon Technologies AG, Neubiberg}
        {Expert Firmware Developer}
        {\begin{itemize}
            \item Entwicklung von Bootstrap Loadern und Startup-Code für verschiedene Mikrocontroller in C.
            \item Aufbau von systematischen Unit-Tests zur Erfüllung von ISO-26262-Anforderungen.
            \item[] \tag{C} \tag{Python} \tag{Ceedling}
        \end{itemize}}
    
    \section[\faMortarBoard]{Studium}
    \job{10/2012 -- 01/2015}
        {Hochschule München, München}
        {Informatik (Master) mit Schwerpunkt Embedded Systems}
        {\begin{itemize}
            \item Masterarbeit bei Fraunhofer ESK: Entwicklung einer Softwarelösung zur hochgenauen Zeit- und Positionsbestimmung in Car2X"=Szenarien
            \item Tätigkeiten in Arbeit: Treiberentwicklung unter GNU/Linux in C
            \item[] \tag{C} \tag{C++}
        \end{itemize}}
    
    \job{10/2008 -- 09/2012}
        {Hochschule München, München}
        {Informatik (Bachelor)}
        {\begin{itemize}
            \item Bachelorarbeit bei Infineon Technologies AG: Entwicklung eines Evaluierungswerkzeugs für barometrische Luftdruck-Sensoren
            \item Tätigkeiten in Arbeit: Firmwareentwicklung für 8- und 16"=Bit"=Mikrocontroller in C und PC"=Software in C\#/.NET
            \item[] \tag{C} \tag{C\#} \tag{.NET}
        \end{itemize}}
    
}
\makebody
\clearpage


\pagestyle{highlightmain}
\highlightbar{

    \section{Interessen}
    \smallskip % additional skip because tag outlines use up space
    % Use \par to group two tags on one line to force a ragged-right layout.
    \tag[black]{Freestyle Football}
    \tag[black]{Fußball}
    \par
    \tag[black]{GNU/Linux}
    \tag[black]{Mnemotechnik}
    \par
    \tag[black]{Klavier spielen}
    \tag[black]{Musik hören}
    \par
    \tag[black]{Rad fahren}
    \tag[black]{Berg gehen}
    
    \section{Was mir wichtig ist}
    \begin{itemize}[itemsep=1mm]
    \item Freude am Lernen: Daher habe ich autodidaktisch das Klavierspielen erlernt.
    \item[] \includegraphics[width=6cm]{Figures/Interests/PianoSheet-59Improvisation-FirstLineCropped}
    \item Ausgleich und Natur: Aus diesem Grund habe ich mit meinem Mountainbike eine Alpenüberquerung gemacht.
    \item[] \includegraphics[width=6cm]{Figures/Interests/CirclesFilledWithNaturePictures}
    \end{itemize}

}
\mainbar{
    
    \section[\faMortarBoard]{Tätigkeiten vor Studium}
    \job{03/2008 -- 09/2008}
    {SEP AG (Entwickler von Datensicherungssoftware), Weyarn}
    {Hilfskraft}
    {\begin{itemize}
            \item Aufsetzen von Testumgebungen (Windows, GNU/Linux), Softwaretests, Pflegen der Anwenderdokumentation.
            \item[] \tag{GNU/Linux} \tag{Windows} \tag{Dokumentation}
    \end{itemize}}

    \job{03/2007 -- 11/2007}
    {Caritas St.-Anna-Haus, Holzkirchen}
    {Zivildienstleistender im Altenheim}
    {\begin{itemize}
            \item Einsatz im Bereich Pflege und Hauswirtschaft.
    \end{itemize}}

    \section[\faMortarBoard]{Schule}
    \job{2004 -- 2006}
        {Staatliche Fachoberschule und Berufsoberschule, Bad Tölz}
        {Fachrichtung: Wirtschaft und Verwaltung}
        {\begin{itemize}
        		\item[] Mit zwei je neunwöchigen Praktika bei Kreissparkasse Miesbach-Tegernsee (Personalabteilung, Hauptbuchhaltung, Kreditabteilung) und Landratsamt Miesbach (Amt für Arbeit und Soziales, Fahrerlaubnisbehörde, Ausländerbehörde)
       	\end{itemize}}
    \job{2000 -- 2004}
        {Staatliche Realschule, Miesbach}
        {}
        {}
    \job{1994 -- 2000}
        {Grund- und Teilhauptschule, Weyarn}
        {}
        {}

    \section[\faFileTextO]{Veröffentlichungen}
    Peer-Reviewed Journalbeiträge:
    \begin{itemize}
        \item \fullcite{kleinmeier-2020}
        \item \fullcite{kleinmeier-2019}
    \end{itemize}

    \vspace{1.5em}

    %Rosenheim, den
    \today\\[0.25em]
    \includegraphics[width=3.5cm]{Figures/Unterschrift/Unterschrift-Kleinmeier-Vectorized}\\[0.25em]
    \textbf{Benedikt Kleinmeier}

}
\makebody

\end{document}
